\chapter*{Avant de commencer}

Ce livre parle d'une machine qui pèse un kilo et demi, consomme vingt watts, comme une ampoule faiblarde, et qui pourtant reconnaît les visages, apprend des langues et attrape une balle au vol. Votre cerveau. Nous allons le regarder comme un ingénieur regarde une carte électronique inconnue: quels sont les composants, comment ils sont reliés, quels signaux courent dans les fils et ce que calcule ce circuit.

Vous n'aurez besoin ni de formules ni de connaissances en biologie. Il vous suffira d'être curieux et de savoir imaginer un fil, une ampoule, un interrupteur. Tout le reste, nous l'assemblerons en chemin.

Ce livre a une seconde version, complète, avec des équations, des modèles, des exercices et des références aux expériences. Les chapitres des deux versions se suivent dans le même ordre et portent les mêmes numéros. Si un chapitre vous accroche, ouvrez sa version complète: vous y trouverez d'où vient chaque affirmation et comment la vérifier vous-même. Sur le site du livre, chaque chapitre passe de \og Court\fg{} à \og Complet\fg{} d'un seul bouton.

Un avertissement. On est loin de tout savoir sur le cerveau, et je dirai honnêtement où s'arrête ce qui est mesuré et où commence la conjecture. Ce n'est pas un défaut du livre, c'est ce qu'il a de plus intéressant: bien des chapitres se terminent sur des questions auxquelles personne ne sait encore répondre.
